\begin{corollary}[Both endpoint conventions are unchanged by rescaling the measure]
Let $\mathcal{M}(\gamma_n;\mathbb{R}^n)$ be the space of $\gamma_n$-a.e.\
equivalence classes of measurable $\mathbb{R}^n$-valued functions on
$\mathbb{R}^n$, and let
$T : L^1(\gamma_n;\mathbb{R}) \to \mathcal{M}(\gamma_n;\mathbb{R}^n)$ be any
linear map with the two properties
\begin{enumerate}
\item $T(\iota v) = R_\gamma v$ for every $v \in L^2(\gamma_n)$, with $\iota$ the
inclusion of \noderef{ToL1} and $R_\gamma$ the transform of
\noderef{RieszTransform}; and
\item $t\,\gamma_n\bigl(\{x : |Tf(x)|_{\mathbb{R}^n} > t\}\bigr)
       \leq 2\,\|f\|_{L^1(\gamma_n)}$
for every $f \in L^1(\gamma_n;\mathbb{R})$ and every real $t > 0$.
\end{enumerate}
Let $c > 0$ be real and let $\mu = c\,\gamma_n$ be the rescaled measure, so
$\mu(A) = c\,\gamma_n(A)$ for every Borel $A$.  Then for every
$f \in L^1(\gamma_n;\mathbb{R})$ and every real $t > 0$, with the \emph{same}
operator $T$:
\begin{enumerate}
\item[(a)] for every $j \in \{1,\ldots,n\}$,
\[
  t\,\mu\bigl(\{x : |(Tf(x))_j| > t\}\bigr) \leq 2 \int |f| \, d\mu ;
\]
\item[(b)]
\[
  t\,\mu\bigl(\{x : |((-\sqrt2)\,T)f(x)|_{\mathbb{R}^n} > t\}\bigr)
    \leq 2\sqrt2 \int |f| \, d\mu .
\]
\end{enumerate}

This is the last sentence of the reference paper's corollary on coordinate and
historical conventions: ``Both statements remain unchanged if $\gamma$ is
replaced by the unnormalized measure $e^{-|x|^2}\,dx$, with the same operator.''
The two statements meant are exactly the two the paper displays --- the
coordinate bound with constant $2$ and the historical bound for
$\mathcal{R} = -\sqrt2\,R_\gamma$ with constant $2\sqrt2$ --- and they reappear
here as (a) and (b).  Since $\gamma_n$ has density $\pi^{-n/2}e^{-|x|^2}$ with
respect to Lebesgue measure (\noderef{GaussMeasure}), the paper's unnormalized
measure $e^{-|x|^2}\,dx$ is the case $c = \pi^{n/2}$; the statement is given for
a general real $c > 0$ because the paper's own justification --- that both the
distribution measure and the input $L^1$ norm acquire the same factor --- is the
general scaling argument and uses nothing about the particular value
$\pi^{n/2}$.

Two points of formulation.  First, the right-hand sides are written as integrals
$\int|f|\,d\mu$ rather than as $\|f\|_{L^1(\mu)}$: this keeps $f$ an element of
$L^1(\gamma_n)$ and $T$ the same operator on the same space, which is what the
paper's ``with the same operator'' asks for, and avoids re-erecting the Lebesgue
space over $\mu$.  This is harmless, since $c > 0$ makes $\mu$ and $\gamma_n$
mutually absolutely continuous, so they have the same null sets, the same
a.e.\ equivalence classes, and the same integrable functions, and
$\int|f|\,d\mu = c\int|f|\,d\gamma_n = c\,\|f\|_{L^1(\gamma_n)}$ is finite.
Second, the content here is neither of the two bounds being rescaled: those hold
over $\gamma_n$, while (a) and (b) are assertions about the measure $\mu$, and
what they record is precisely that the weak-type constants $2$ and $2\sqrt2$ are
invariant under the rescaling, i.e.\ that they are properties of the operator and
not artifacts of the normalization of the underlying measure.

Properties (1) and (2) are written out here rather than imported by citation.
They are exactly the two properties for which \noderef{EndpointTheorem} asserts
existence and uniqueness of a map $T$, so the universally quantified statement
above is not vacuous and its instance at that unique $T$ is the paper's
assertion.
\end{corollary}
\begin{proof}
Fix $c > 0$, $f \in L^1(\gamma_n;\mathbb{R})$ and a real $t > 0$, and put
$\mu = c\,\gamma_n$.

\emph{Step 1: both sides carry the same factor $c$.}
By definition of the rescaled measure, $\mu(A) = c\,\gamma_n(A)$ for every Borel
set $A$, as an identity in $[0,+\infty]$.  Likewise, integration against
$c\,\gamma_n$ is $c$ times integration against $\gamma_n$, so
\begin{equation*}
  \int |f| \, d\mu = c \int |f| \, d\gamma_n = c\,\|f\|_{L^1(\gamma_n)} ,
  \tag{1}
\end{equation*}
the second equality being the definition of the $L^1(\gamma_n)$ norm of $f$ as
the integral of $|f|$ against $\gamma_n$; both quantities are finite because
$f \in L^1(\gamma_n)$.

\emph{Step 2: part (a).}
Let $j \in \{1,\ldots,n\}$.  By \noderef{CoordinateEndpointBound}, applied to
this $T$ with its hypotheses (1) and (2) --- which are the hypotheses of the
present statement verbatim --- and to this $f$, this $j$ and this $t > 0$,
\[
  t\,\gamma_n\bigl(\{|(Tf)_j| > t\}\bigr) \leq 2\,\|f\|_{L^1(\gamma_n)} .
\]
Multiplying this inequality in $[0,+\infty]$ by the nonnegative constant $c$,
which preserves the order, and using Step 1 on both sides,
\[
  t\,\mu\bigl(\{|(Tf)_j| > t\}\bigr)
    = c\;t\,\gamma_n\bigl(\{|(Tf)_j| > t\}\bigr)
    \leq c \cdot 2\,\|f\|_{L^1(\gamma_n)}
    = 2 \int |f| \, d\mu ,
\]
which is (a).

\emph{Step 3: part (b).}
By \noderef{SpectralNormalizationBound}, applied to the same $T$ with the same
two hypotheses and to this $f$ and this $t > 0$,
\[
  t\,\gamma_n\bigl(\{|((-\sqrt2)T)f|_{\mathbb{R}^n} > t\}\bigr)
    \leq 2\sqrt2\,\|f\|_{L^1(\gamma_n)} .
\]
Multiplying by $c \geq 0$ and using Step 1 on both sides exactly as in Step 2,
\[
  t\,\mu\bigl(\{|((-\sqrt2)T)f|_{\mathbb{R}^n} > t\}\bigr)
    = c\;t\,\gamma_n\bigl(\{|((-\sqrt2)T)f|_{\mathbb{R}^n} > t\}\bigr)
    \leq c \cdot 2\sqrt2\,\|f\|_{L^1(\gamma_n)}
    = 2\sqrt2 \int |f| \, d\mu ,
\]
which is (b).  In both parts the factor $c$ appears once on each side and
cancels, which is the reference paper's ``scaling the measure scales both sides
by the same factor''.
\end{proof}
